\section{Remaining chromatic and structural questions}
\label{sec:higher}
The all-dimensional clique number is settled, and the exact dimension-six
chromatic numbers are $\chi(\Gamma_6)=12$ and $\chi(\OG_6^*)=15$.
In dimension seven the line and full-subspace graphs require at least
nineteen colors. The clique number of the full graph is sixteen, so the
chromatic--clique separation is at least three.

The exact dimension-seven chromatic numbers remain open. No proper
nineteen-coloring of either full graph is supplied here. The proof of the
lower bound in Section~\ref{sec:seven} excludes eighteen colors; it neither
excludes nineteen nor relies on a complete classification of
nineteen-color profiles. The partial count, charge, orbit and covering
analysis of such profiles is retained in the separate extended working
draft and the repository's proof notes and certificates.

For $n\geq8$, determining the chromatic number remains open, as does the
question whether $\chi(\OG_n^*)>\omega(\OG_n^*)$ holds in infinitely many
dimensions. These questions require further chromatic arguments; the
clique formula alone does not settle them.

Within dimension six, a structural description of an optimal full-subspace
coloring remains open. Appendix~\ref{sec:geometry} retains the established
obstructions: exact profiles and the specified lift features do not
separate all adjacent vertices, and equivariance under the full
stabilizer of the fixed totally isotropic three-space is impossible with
the prescribed clique labels. Auxiliary geometric data or a smaller
symmetry group may provide a conceptual construction.

\begin{problem}\label{prob:structural}
Can one construct a proper fifteen-coloring of $\OG_6^*$ from explicit
geometric data, proving availability and separation without looking up
the archived coloring certificate?
\end{problem}
The checked twelve-coloring of $\Gamma_6$ establishes its chromatic
number, but full orthogonal-group equivariance is not established here.
